\section{Combining the two forests, and a linear-algebra extension}
\label{sec:combination}

\subsection{An exact coupled finite problem}

Suppose an authenticated handle model supplies a primal graph $G=(V,E)$,
a connected rooted dual graph $D$, and an occurrence matrix
$N\in\NN^{R\times(E\sqcup J)}$.  Each row restricts to a closed walk in
$G$.  Eligible rows are in bijection with the dual edges; other rows may
be ineligible for top-handle deletion.  Let $\mathcal L$ be the total bit
length of the graphs, matrix, and row/edge correspondence.
Assume every dual spanning tree is a legal specified
two/three-handle reduction that deletes precisely its selected rows,
leaves the retained occurrence data unchanged, and preserves the exposed
ports required by the subsequent operation.  These are geometric
hypotheses on the model, not automatic properties of a pair of graphs.
Also assume that every chosen primal spanning forest has the authenticated
forest-first contraction used in Section~\ref{sec:deferred-primal-forest}:
it deletes those columns and preserves the retained attaching charts and
required port interpretation.  This lower contraction contract is part
of the coupled geometric application.

\begin{corollary}[Coupled forests without primal enumeration]
\label{cor:coupled-forests}
Under these hypotheses, there exist a legal dual tree and a primal
spanning forest permitting complete raw positive cancellation if and
only if there exists a dual tree whose deleted-row complement is
completely column-peelable on \emph{all} columns $E\sqcup J$.
The latter problem is decidable in
\[
 2^{O(u\log(u+2))}\operatorname{poly}(\mathcal L),
\]
where $u$ is the exceptional-interface size of the full pre-deletion
matrix.  A successful decision constructs both trees and a complete
raw cancellation sequence.
\end{corollary}
\begin{proof}
Fix a legal dual tree $T_D$.  Every retained row still has its original
closed-walk interpretation, so Theorem~\ref{thm:all-primal-forests} applies
to the row restriction $R\setminus T_D$.  It gives the stated equivalence
for this particular tree.  Quantifying over the dual trees proves the
first assertion.  Apply Theorem~\ref{thm:exceptional-selector} to the full
matrix and dual graph; after its successful tree is found, peel the
retained rows and choose a primal forest avoiding the graph pivots.
The extra finite work is polynomial and introduces no new search factor.
\end{proof}

For original native source-core columns only, the prior bound $u\le3k$ applies when
$k$ counts exceptional columns in this full input.  It must not be
estimated only after a favorable but unknown dual tree is chosen: that
would assume the answer in the parameter bound.  A separate parameterized
algorithm would be needed if the parameter were instead the minimum
exceptional size over all dual choices.

The delivered native patch uses the existing breadth-first dual policy
and the new primal theorem.  The delivered exceptional solver is a
separate finite operator.  Corollary~\ref{cor:coupled-forests} justifies a
future source adapter under the stated contracts; it is not a claim that
the current whole-cut wrapper already implements exhaustive coupled
selection.  In particular, no extra native successes are attributed to
the standalone dual solver in our result table.

\subsection{Deferred bases in a represented binary matroid}

The connectivity argument is a special case of a useful algebraic fact.
It also identifies exactly what would be lost by replacing a forest with
a basis of a represented binary matroid.

\begin{theorem}[Deferred column basis]
\label{thm:deferred-basis}
Let $B$ be a matrix over $\mathbb F_2$ with columns indexed by $E$, and
let $N\in\NN^{R\times(E\sqcup J)}$.  Assume that each uncapped row
vector on $E$ lies modulo two in the kernel of $B$:
\[
 B\,(n_{re}\bmod2)_{e\in E}=0\qquad(r\in R).
\]
If a column-leaf sequence in $N$ uses $E$-indexed pivot set
$P\subseteq E$, then
\[
 \operatorname{span}(B_{E\setminus P})=
 \operatorname{span}(B_E).
\]
Consequently $E\setminus P$ contains a column basis $K$ of $B$.
Moreover, full-matrix peeling succeeds if and only if some column basis
$K$ can be discarded while leaving a matrix admitting complete raw
positive singleton cancellation.
\end{theorem}
\begin{proof}
Before a pivot $e\in E$ in row $r$, every earlier pivot column has
coefficient zero in $r$.  Its coefficient at $e$ is one.  The kernel
identity expresses $B_e$ as a sum of the other currently retained columns
with coefficients $n_{rf}\bmod2$.  Deleting $e$ therefore preserves
the span.  Auxiliary pivots change none of the columns of $B$.
Induction proves the displayed equality, so a basis can be selected from
the complement of all pivots.

If full peeling succeeds, choose this basis $K$.  The recorded sequence
uses no column of $K$ and remains valid after those columns are discarded.
Conversely, complete raw positive elimination after a basis deletion
implies complete peeling by Theorem~\ref{thm:raw-completeness}, and the
same pivot-column sums make that peeling sequence valid with the deleted
basis columns restored.  This proves the equivalence.
\end{proof}

Take $B$ to be the mod-two vertex/edge incidence matrix of a graph.  Its
column bases are exactly spanning forests, and the theorem recovers the
primal result.  More generally, a cellular boundary relation
$\partial_k\partial_{k+1}=0$ can supply the kernel hypothesis for a
represented matroid, provided that the raw incidences have the claimed
mod-two interpretation.  This gives an algebraic normalization principle
in other chain complexes.

An arbitrary column basis in a higher boundary matrix need not be
geometrically collapsible.  Therefore this generalization does not
authorize higher-dimensional handle cancellation, replacement of a
marked disk system, or compression of a manifold from homology alone.
The special advantage of dimension one is that a graph basis is an
forest with contractible components.  Proving a similarly useful geometric
realization criterion in another setting is a research problem.

\subsection{Stopping sets as reusable restricted certificates}

The universal primal obstruction is unusually cheap to store: a nonempty
list of source row identifiers suffices.  Source replay checks that every
column total on those rows avoids one.  Such a witness can be reused
across all primal forest policies on the same retained rows and charts.
It invalidates neither a different dual deletion nor a different
geometric presentation.  A cache must therefore bind the witness to the
authenticated source, retained-row set, matrix version, and occurrence
convention.  Using an obstruction from an earlier hierarchy state after
changing those data would be unsound.

The exceptional solver's exhaustive negative has a different status.
Its search proves failure for every allowed dual tree in one finite
instance, but the current implementation provides no compact
independently replayable certificate of the entire exhaustion.  These
two kinds of refusal should have separate names in an eventual API.
